\documentclass{TexHU}
\usepackage{lipsum}
\usepackage{pgfplots}
\usepackage{adjustbox}
\usepackage{titlesec}
\usepackage{xcolor}
\usepackage{ulem}
\pgfplotsset{compat=1.18}

\title{\LaTeX\,Test for libertine 中文测试}
\author{Tester}

% \newcommand{\sectiondivider}{%
%   \vspace{0.1in}\noindent\makebox[\linewidth]{\dashuline{\hspace{\linewidth}}}%
% }
\newcommand{\sectiondivider}{%
  \vspace{0.1in}
  \centerline{\dashuline{\hspace{0.4\linewidth}}}
}

\titleformat{\section}
  {\normalfont\Large\bfseries}
  {}
  {0pt}
  {\color{black}\rule{\dimexpr\fontcharht\font`X}{\dimexpr\fontcharht\font`X}\quad}


\begin{document}
% \maketitle
\begin{adjustbox}{width=\textwidth}
    \begin{minipage}[c]{\linewidth}
        \begin{tikzpicture}[remember picture, overlay]
            \draw[line width=0.1in] (0in,-0.5in) -- ++(0, -1in);
            \node [anchor=north west, align=left, minimum height=1in] at ([xshift=1.25in, yshift=-1.5in]current page.north west) {%
                \LARGE\textbf{现代计算机体系架构：综合课程设计}\\
                \large\bfseries\textsc{Modern Computer Architecture: Final Project Design}
            };
            \node [anchor=north east] at ([xshift=-0.5in, yshift=-0.5in]current page.north east) {%
                \includegraphics[height=0.95cm]{assets/thu.eps}
            };
            \node [anchor=north west, align=left] at ([xshift=0.5in, yshift=-0.5in]current page.north west) {%
            \footnotesize \cfs 姓名：刘国鸿 \\
            \footnotesize \cfs 学号：{\cfsEn 2020010688}
            };
            % \node [anchor=north] at ([yshift=-2.75in]current page.north) {%
            % {\small\cfs 姓名：刘国鸿\hspace{2em}学号：{\cfsEn 2020010688}\hspace{2em}日期：\cfsEn{2023/06/13}}
            % };
        \end{tikzpicture}
        \vspace{2in}
    \end{minipage}
\end{adjustbox}

\section*{实验介绍}

神经网络是目前最热的技术之一。神经网络加速器是为了提高神经网络算法而提出的新型计算单元。与传统的CPU、GPU相比，神经网络加速器可以大幅度提高性能、降低功耗。然而，高效的神经网络加速器通常只支持有限的网络算子，而在完整的神经网络应用中，通常还需要一些其他加速设备对数据进行预处理等其他操作。

本次大作业提供了一个简单的神经网络加速器仿真框架，通过实验研究如何进行涉及多个不同加速器的系统仿真，并尝试一些优化方案。


This is a \textbf{random} text for \textsc{testing} the \textit{performance} of such a new font.

Let's try some different ligatures, such as !!, !?, \textit{ffl}, ffi.

There is roughly a \(50\%\) 中文插入chance.

中文测试。具体测试包括\textbf{粗体}，\textit{斜体}，等等。

\cfs 中文测试 testing for English.
\[a^2+b^2=c^2=\pi\sqrt[3]{a}\]

\sectiondivider

\section*{实验内容}



\end{document}